% AI-written, AI cross-reviewed; no human sign-off; not a claim that RH is proved
% Atlas unit P025. Local manuscript v3, item 7.2, heading on printed PDF p. 72.
% UpperTheory.V.Th.adjointLiftingTheorem_of_locPrints_of_core
% paper2/inputs_one.tex:46–80. Verbatim excerpt; SHA-256 80253032c1c655c9fc12089c5d959573b2397a645e55d8f612b63b38fba8ab48.
% Original TeX uses the macros and cross-reference keys of the source manuscript.


% Start of atlas edition annotation. The following block is outside the original excerpt.
\paragraph{Status of named inputs in this reading.}
\noindent\nolinkurl{UpperTheory.V.Loc.LocalPublishedTheorems.bk} --- record field (\path{Development/AutomorphicArithmetic/AutomorphicVarieties/PadicHodgeTheory/LocalPublishedTheorems.lean:19}); not compared.\par
\noindent\nolinkurl{UpperTheory.V.Loc.LocalPublishedTheorems.lfK} --- record field (\path{Development/AutomorphicArithmetic/AutomorphicVarieties/PadicHodgeTheory/LocalPublishedTheorems.lean:13}); not compared.\par
\noindent\nolinkurl{UpperTheory.V.Loc.LocalPublishedTheorems.ltK} --- record field (\path{Development/AutomorphicArithmetic/AutomorphicVarieties/PadicHodgeTheory/LocalPublishedTheorems.lean:14}); not compared.\par
\noindent\nolinkurl{UpperTheory.V.Loc.LocalPublishedTheorems.leK} --- record field (\path{Development/AutomorphicArithmetic/AutomorphicVarieties/PadicHodgeTheory/LocalPublishedTheorems.lean:15}); not compared.\par
\noindent\nolinkurl{UpperTheory.V.Loc.LocalPublishedTheorems.lf} --- record field (\path{Development/AutomorphicArithmetic/AutomorphicVarieties/PadicHodgeTheory/LocalPublishedTheorems.lean:16}); not compared.\par
\noindent\nolinkurl{UpperTheory.V.Loc.LocalPublishedTheorems.lt} --- record field (\path{Development/AutomorphicArithmetic/AutomorphicVarieties/PadicHodgeTheory/LocalPublishedTheorems.lean:17}); not compared.\par
\noindent\nolinkurl{UpperTheory.V.Loc.LocalPublishedTheorems.le} --- record field (\path{Development/AutomorphicArithmetic/AutomorphicVarieties/PadicHodgeTheory/LocalPublishedTheorems.lean:18}); not compared.\par
% End of atlas edition annotation.

\begin{formalstatement}[Ordinary to Bloch--Kato adjoint vanishing]\label{pn:P025}
Assume local cohomological finiteness, Tate duality, the Euler
characteristic formulas and the local Bloch--Kato theorem in the
specified coefficient models. Assume the dense ordinary vanishing
hypothesis: at every locally generic automorphic point below each
allowed ordinary adjoint cocycle is a coboundary. Let \(L\) be CM with
\(\mu_\ell\not\subset L\), take a Hida setting and level \(V\) with
\(V.S'=S\), large coefficients, \(2\le n\le4\), an admitted residual
representation, and a locally generic automorphic point \(x,\rho\)
of its universal ordinary determinant ring. For each \(v\mid\ell\)
suppose there are strictly decreasing integers \(m_i\) and \(k>0\)
such that on inertia
\[
 \bigl(\psi_{v,i}\psi_{v,j}^{-1}\chi^{m_j-m_i}\bigr)^k=1
 \quad(i,j=1,\ldots,n).
\]
For every allowed polarizing form and adjoint extension \(R\), each
continuous cocycle \(z\) whose pullback satisfies the Bloch--Kato
finite condition has the form \(z(g)=R(g)X-X\) for one matrix \(X\)
and all \(g\).
\end{formalstatement}\nodestatus{P025}

The scalar-tower identity identifies the two multiplier characters in
the ordinary and finite conditions. The local comparison must transport
the period primitive, descend it on the stable corner, and recover the
inertia flag. The global ordinary vanishing statement is a premise.
Newton and Thorne's \emph{Adjoint Selmer groups of automorphic Galois
representations of unitary type}, Theorem~A, uses an enormous-image
hypothesis over the cyclotomic extension. The listed assumptions do
not supply that match. P.~B.~Allen's \emph{Deformations of polarized
automorphic Galois representations and adjoint Selmer groups}, Theorem~A,
requires adequate residual image after restriction. Blueprint Z1.2~\bploc{r6:review:Z1.2}
records the comparison from the project finite condition to its ordinary
flag condition separately from the production of ordinary vanishing.
\cite{r21-nt}\quad\cite{r21-allen}
